% =============================================================================
%  QED-SAPT0 in the cqed-scf package: the density-fitted Pauli-Fierz integral
%
%  Overleaf-ready: a single self-contained file, no external .bib required.
%
%  Companion documents:
%    docs/qed_cis_formalism.tex                      -- the separable-kernel
%                                                       observation this note
%                                                       extends (its Sec. 2.3)
%    docs/development/QED_SAPT_DISPERSION_PLAN.md    -- implementation plan
%    docs/development/psi4_array_aliasing.md         -- the aliasing defect note
%    docs/SAPT_NOTE.md                               -- running development log
% =============================================================================
\documentclass[11pt]{article}

\usepackage[margin=1in]{geometry}
\usepackage{amsmath, amssymb, mathtools, bm}
\usepackage{booktabs}
\usepackage{braket}
\usepackage{enumitem}
\usepackage{microtype}
\usepackage[dvipsnames]{xcolor}
\usepackage[colorlinks=true, linkcolor=BrickRed, citecolor=BrickRed,
            urlcolor=NavyBlue]{hyperref}

% ---- macros (kept consistent with qed_cis_formalism.tex) --------------------
\newcommand{\lam}{\bm{\lambda}}
\newcommand{\dop}{\hat{d}}                      % lambda . mu_el (one-electron)
\newcommand{\dexp}{\langle d \rangle}
\newcommand{\Hpf}{\hat{H}_{\mathrm{PF}}}
\newcommand{\naux}{N_{\mathrm{aux}}}
\newcommand{\nbf}{N_{\mathrm{bf}}}
\DeclareMathOperator{\Tr}{Tr}

\title{\textbf{QED-SAPT0 in \texttt{cqed-scf}}\\[2pt]
\large The dipole self-energy as a single density-fitting auxiliary function}
\author{}
\date{}

\begin{document}
\maketitle
\vspace{-3em}

\begin{abstract}
\noindent
This note documents the integral factorization underlying the density-fitted
QED-SAPT0 backend in \texttt{cqed-scf}. Its companion,
\texttt{docs/qed\_cis\_formalism.tex}, records the structural observation that
the two-electron dipole self-energy (DSE) is an electron-repulsion integral
with a rank-one separable kernel. Read in the language of density fitting, that
observation says something stronger and considerably more useful: the DSE is
\emph{exactly one additional auxiliary basis function}. Appending a single row
to the density-fitting three-index tensor therefore converts every
standard-SAPT working equation into its Pauli--Fierz counterpart without
altering a single formula, sign convention, or contraction pattern. We derive
the result, state the symmetry condition it requires and why that condition
holds structurally, give the resulting expressions for the second-order
dispersion and exchange-dispersion terms, and record the numerical evidence.
The practical consequence is a sharp diagnostic: the cavity contribution
carries \emph{no} fitting error, so cavity blocks must agree with a dense
reference to machine precision while standard blocks agree only to
density-fitting error.
\end{abstract}

\tableofcontents
\bigskip

% =============================================================================
\section{Scope and notation}
\label{sec:scope}

We treat closed-shell QED-SAPT0 with monomer CQED-SCF references in the
dimer-centred basis (the ``DCBS'' convention: each monomer is computed with the
partner's basis functions present as ghosts). Orbital-index conventions follow
the driver implementation and \texttt{psi4numpy}'s \texttt{helper\_SAPT}:

\begin{center}
\small
\begin{tabular}{ll}
\toprule
symbol & meaning \\
\midrule
$a, a'$ & occupied orbitals of monomer $A$ \\
$r, r'$ & virtual orbitals of monomer $A$ \\
$b, b'$ & occupied orbitals of monomer $B$ \\
$s, s'$ & virtual orbitals of monomer $B$ \\
$p,q,\mu,\nu$ & atomic-orbital (AO) indices in the shared dimer basis \\
$P, Q$ & auxiliary (fitting) basis indices \\
\bottomrule
\end{tabular}
\end{center}

\noindent
$\dop = \lam\cdot\hat{\bm{\mu}}_{\mathrm{el}}$ is the dipole operator projected
onto the cavity polarization, with AO matrix
\begin{equation}
d_{\mu\nu} = \sum_{x \in \{x,y,z\}} \lambda_x \, \mu^{x}_{\mu\nu} .
\label{eq:dao}
\end{equation}
We write $\dexp_X$ for the electronic coherent-state expectation value over
monomer $X$'s CQED-SCF reference.

% =============================================================================
\section{The Pauli--Fierz two-electron integral}
\label{sec:pf}

\subsection{The separable kernel, restated}

The companion document establishes that the two-electron part of the DSE,
$\tfrac{1}{2}\sum_{i\neq j}\dop(i)\dop(j)$, is a two-electron operator whose
kernel is the rank-one product
\begin{equation}
(pq|rs) \;\longrightarrow\; d_{pq}\, d_{rs} ,
\label{eq:sepkernel}
\end{equation}
so that every Slater--Condon result carries over verbatim under that
substitution. For the intermolecular SAPT problem the relevant object is the
$A$--$B$ cross term of the coherent-state fluctuation product, which supplies
the \emph{Pauli--Fierz} two-electron integral
\begin{equation}
(pq|rs)_{\mathrm{PF}} \;=\; (pq|rs) \;+\; d^{A}_{pq}\, d^{B}_{rs} ,
\label{eq:pfint}
\end{equation}
where $d^{A}$ and $d^{B}$ are the dipole-projected AO matrices built in
monomer $A$'s and monomer $B$'s (ghosted) bases.

\paragraph{Coefficient.} In the coherent-state basis the DSE is
$\tfrac{1}{2}\bigl(\lam\cdot(\hat{\bm\mu}-\langle\hat{\bm\mu}\rangle)\bigr)^2$,
so the $A$--$B$ cross term is
$(\lam\cdot\Delta\bm{\mu}_A)(\lam\cdot\Delta\bm{\mu}_B)$ with coefficient
unity: the $\tfrac12$ cancels against the cross term appearing twice. This is
the coefficient in Eq.~\eqref{eq:pfint}.

\paragraph{Why the $\dexp$ shift does not appear here.} Subtracting a constant
times the identity shifts only \emph{diagonal} matrix elements, so the
off-diagonal (transition) blocks that enter the dispersion numerator are
untouched. The shift does matter for the first-order and induction terms, where
it contributes the one-electron potentials
$V^{A,\mathrm{cav}} = -\dexp_A\, d^{B}$ and
$V^{B,\mathrm{cav}} = -\dexp_B\, d^{A}$ and the scalar $\dexp_A \dexp_B$.
Those are $O(N^2)$ objects handled outside the factorization developed here.

\subsection{The symmetry condition}
\label{sec:symmetry}

Equation~\eqref{eq:pfint} is written with distinct labels $d^A$ and $d^B$, but
by Eq.~\eqref{eq:dao} the matrix $d$ depends only on $\lam$ and on the AO
basis and frame in which the dipole integrals are evaluated --- \emph{not} on
which monomer's electrons occupy it. In the DCBS convention both monomers are
computed in the same dimer basis and the same frame, hence
\begin{equation}
d^{A}_{\mu\nu} = d^{B}_{\mu\nu} \;\equiv\; d_{\mu\nu} ,
\qquad d = d^{\mathsf T} .
\label{eq:dAdB}
\end{equation}
This has been verified numerically to be an exact equality (bitwise
$\max|d^A - d^B| = 0$), both with and without \texttt{no\_com}/%
\texttt{no\_reorient}. It is asserted in the code rather than assumed: a
violation would mean the two monomer bases or dipole origins had diverged,
which invalidates far more than the factorization below.

With Eq.~\eqref{eq:dAdB}, the Pauli--Fierz kernel
\begin{equation}
(pq|rs)_{\mathrm{PF}} = (pq|rs) + d_{pq}\, d_{rs}
\label{eq:pfsym}
\end{equation}
is symmetric under $(pq)\leftrightarrow(rs)$ and positive semidefinite in the
added term --- exactly the structural properties a density-fitting expansion
possesses.

% =============================================================================
\section{The DSE as one auxiliary function}
\label{sec:onerow}

\subsection{Derivation}

Ordinary robust density fitting expands the Coulomb kernel in an auxiliary
basis of $\naux$ functions,
\begin{equation}
(pq|rs) \;\approx\; \sum_{PQ} (pq|P)\,[\bm{J}^{-1}]_{PQ}\,(Q|rs)
\;=\; \sum_{P=1}^{\naux} B^{P}_{pq}\, B^{P}_{rs} ,
\qquad
B^{P}_{pq} \equiv \sum_{Q} [\bm{J}^{-1/2}]_{PQ}\,(Q|pq),
\label{eq:df}
\end{equation}
with metric $J_{PQ} = (P|Q)$. Equation~\eqref{eq:df} is the statement that the
Coulomb kernel is approximated by a sum of $\naux$ separable rank-one terms.

The DSE kernel of Eq.~\eqref{eq:pfsym} \emph{is} a separable rank-one term, of
precisely the same algebraic form, with $d$ in place of $B^P$. Therefore
\begin{equation}
\boxed{\;
(pq|rs)_{\mathrm{PF}} \;\approx\; \sum_{P=1}^{\naux+1} \tilde{B}^{P}_{pq}\,\tilde{B}^{P}_{rs},
\qquad
\tilde{B}^{P} =
\begin{cases}
B^{P}, & P = 1,\dots,\naux \quad\text{(Coulomb, fitted)}\\[3pt]
d, & P = \naux+1 \quad\text{(DSE, exact)}
\end{cases}
\;}
\label{eq:augmented}
\end{equation}

\noindent
\textbf{The dipole self-energy is one extra auxiliary basis function.} No
approximation has been introduced by the augmentation itself: rows
$1,\dots,\naux$ carry the usual fitting error, and row $\naux+1$ is exact.

\subsection{Three consequences}

\begin{enumerate}[leftmargin=*, itemsep=4pt]

\item \textbf{Every working equation is unchanged.} Any expression written in
terms of $B$ becomes its Pauli--Fierz counterpart by substituting $\tilde{B}$
and extending the auxiliary summation by one. There is no parallel ``cavity''
code path, no additional sign convention to reconcile, and no new intermediate.
Given the number of sign-convention traps already documented for this package
(Psi4's \texttt{Hx}-family routines each combine and sign their results
differently), eliminating a code path is worth more than the arithmetic it
saves.

\item \textbf{Operator contexts are slices of the auxiliary index.} The
partition used throughout the driver,
\begin{equation}
\underbrace{P \le \naux}_{\textsf{standard}}, \qquad
\underbrace{P = \naux+1}_{\textsf{cavity}}, \qquad
\underbrace{P \le \naux+1}_{\textsf{total}},
\label{eq:contexts}
\end{equation}
is exact and free: \textsf{standard} $+$ \textsf{cavity} $=$ \textsf{total}
holds to roundoff by construction, not by cancellation of separately computed
quantities. Disabling the cavity is simply omitting the last row, whereupon the
\textsf{cavity} slice is empty and contracts to exact zeros with no special
case.

\item \textbf{The cavity terms carry no fitting error, which is a sharp test.}
Because row $\naux+1$ is exact, any quantity built purely from the
\textsf{cavity} context must reproduce a dense reference to \emph{machine
precision}, while \textsf{standard} quantities agree only to density-fitting
error. An implementation error in the pairing, transposition, or scaling of the
DSE row would degrade cavity agreement to fitting-error level --- a five-order-%
of-magnitude signal. See Sec.~\ref{sec:numerics}.

\end{enumerate}

\subsection{Cost}

The augmentation changes $\naux \to \naux + 1$; every scaling is unchanged.
The saving relative to the dense four-index formulation is the point. For the
water/methylamine reference dimer:

\begin{center}
\small
\begin{tabular}{lrrrr}
\toprule
basis & $\nbf$ & $\naux$ & dense $\nbf^4$ (3 tensors) & DF $\naux\,\nbf^2$ \\
\midrule
jun-cc-pVDZ & 89 & 293 & 1.51\,GB & 18\,MB \\
aug-cc-pVDZ & 132 & 377 & 7.29\,GB & 52\,MB \\
aug-cc-pVTZ & 299 & 640 & 192\,GB & 458\,MB \\
\bottomrule
\end{tabular}
\end{center}

\noindent
The dense path stores three $\nbf^4$ tensors (standard, cavity, total). Note
that the cavity tensor is rank one: materializing $d\otimes d$ spends $\nbf^4$
doubles to represent $\nbf^2$ numbers. Equation~\eqref{eq:augmented} never
forms it.

% =============================================================================
\section{The pair factorization of the SAPT integral blocks}
\label{sec:pairs}

\subsection{The pairing rule}

SAPT working equations are written in terms of four-index blocks that the
driver addresses by a four-character string. The internal storage convention is
\begin{equation}
\mathcal{I}[p,q,r,s] \;=\; (p\,r\,|\,q\,s)_{\mathrm{PF}}
\end{equation}
(chemists' notation with the middle indices transposed), and the block
$v(\sigma_0\sigma_1\sigma_2\sigma_3)$ is obtained by transforming index $k$ with
the orbital coefficients of space $\sigma_k$. Inserting
Eq.~\eqref{eq:augmented} and collecting terms gives the central result of this
section:
\begin{equation}
\boxed{\;
v(\sigma_0\sigma_1\sigma_2\sigma_3)_{ABCD}
\;=\; \sum_{P=1}^{\naux+1}
b^{P}_{(\sigma_0\sigma_2)}[A,C]\;
b^{P}_{(\sigma_1\sigma_3)}[B,D] \;} ,
\label{eq:pairfac}
\end{equation}
where the \emph{pair block} is the half-transformed three-index tensor
\begin{equation}
b^{P}_{(\sigma\tau)}[A,C] \;=\; \sum_{\mu\nu} C^{\sigma}_{\mu A}\,
\tilde{B}^{P}_{\mu\nu}\, C^{\tau}_{\nu C} .
\label{eq:pairblock}
\end{equation}

\noindent
\textbf{The pairing is positions $(0,2)$ and $(1,3)$}, not $(0,1)$ and $(2,3)$:
the first pair is the monomer-$A$ pair and the second the monomer-$B$ pair.
This is the algebraic content of the string convention, and it is what makes a
\emph{two}-character primitive the natural one --- the same arity as the
overlap $s(\sigma\tau)$ and one-electron potential blocks the driver already
exposes. The four-index $v$ is then derived, not primitive.

\paragraph{Cross-monomer pairs are required.} The seventeen strings appearing
in the driver need ten distinct pairs,
\[
aa,\; ab,\; ar,\; as,\; ra,\; ba,\; bb,\; br,\; bs,\; sb ,
\]
of which $ab$, $as$, $br$, $sb$, $ba$ and $ra$ mix the two monomers. A design
that only formed same-monomer pair blocks would not support
exchange-dispersion.

\paragraph{Storage.} Because all four orbital spaces live in the same dimer AO
basis, half-transforming $\tilde{B}$ once into the union space
$C_{\mathrm{all}} = [\,C_a \,|\, C_r \,|\, C_b \,|\, C_s\,]$ makes every pair
block a \emph{slice} of a single $(\naux+1)\times 2N_{\mathrm{mo}} \times
2N_{\mathrm{mo}}$ array. All sixteen pairs are then available at the cost of
the largest one, and pair lookup costs no arithmetic.

\paragraph{Implementation note: the fitting basis must be chosen by role.}
The auxiliary expansion of Eq.~\eqref{eq:df} is not one fixed object. Fitting
sets are optimized for a purpose: JKFIT sets for Coulomb and exchange matrices,
RIFIT sets for correlation energies. The distinction is quantitatively decisive
for SAPT because $E^{(10)}_{\mathrm{elst}}$ is a small residual of large
cancelling contributions --- the two-electron Coulomb term against nuclear
attraction and nuclear repulsion --- so fitting error in the two-electron part
cancels against nothing and survives at full size in the residual. Correlation
energies have no such cancellation.

Measured on water/He at cc-pVDZ, using a RIFIT set for every term gives an
$E^{(10)}_{\mathrm{elst}}$ error of $3.2\times10^{-5}\,E_h$ against
$4.0\times10^{-7}\,E_h$ for JKFIT, a factor of $79$; the exchange-induction and
induction terms improve by factors of $53$ and $19$. Dispersion is indifferent
(factor $0.9$), being the term RIFIT was designed for. This mirrors Psi4's own
SAPT, which fits electrostatics, exchange and induction with
\texttt{DF\_BASIS\_ELST}/\texttt{DF\_BASIS\_SCF} (JKFIT) and the MP2-like terms
with \texttt{DF\_BASIS\_SAPT} (RIFIT).

The implementation therefore maintains two augmented tensors, differing only in
their first $\naux$ rows; the DSE row $\tilde B^{\naux+1} = d$ is identical in
both and exact in both. The choice of fitting set is an approximation to the
Coulomb kernel alone and never touches the cavity contribution.

\paragraph{Implementation note (aliasing).} The three operator contexts of
Eq.~\eqref{eq:contexts} and all pair blocks are overlapping views of that one
array. They are marked read-only. An in-place write would silently couple the
contexts --- the same class of defect recorded in
\texttt{docs/development/psi4\_array\_aliasing.md}, where a NumPy view of a
Psi4 buffer collapsed two physically distinct orbital-energy sets into one and
produced plausible, wrong numbers for an extended period.

\subsection{Dressed blocks}

The SAPT interaction operator $\tilde{v}$ adds one-electron and scalar
dressings to the two-electron block,
\begin{equation}
\tilde{v}(\sigma_0\sigma_1\sigma_2\sigma_3)_{ABCD}
= v_{ABCD}
+ S^{(\sigma_0\sigma_2)}_{AC}\,\frac{V^{A,(\sigma_1\sigma_3)}_{BD}}{N_A}
+ \frac{V^{B,(\sigma_0\sigma_2)}_{AC}}{N_B}\,S^{(\sigma_1\sigma_3)}_{BD}
+ S^{(\sigma_0\sigma_2)}_{AC} S^{(\sigma_1\sigma_3)}_{BD}\,
  \frac{\mathcal{V}_{\mathrm{nuc}}}{N_A N_B},
\label{eq:vt}
\end{equation}
with $N_X$ the electron count of monomer $X$. Every term on the right
factorizes over the \emph{same} pair partition $(\sigma_0\sigma_2)$,
$(\sigma_1\sigma_3)$ as Eq.~\eqref{eq:pairfac}. The cavity contributions to
$V^A$, $V^B$ and $\mathcal{V}_{\mathrm{nuc}}$ are the $\dexp$-dependent pieces
recorded in Sec.~\ref{sec:pf}; all are $O(N^2)$ or $O(N^3)$ and require no
fitting. Consequently \textbf{only $v$ needs re-plumbing}: $\tilde v$ inherits
the factorization unchanged.

% =============================================================================
\section{Second-order dispersion}
\label{sec:disp}

\subsection{Amplitudes and \texorpdfstring{$E^{(20)}_{\mathrm{disp}}$}{Edisp20}}

The uncoupled dispersion amplitudes and energy are
\begin{align}
t_{rsab} &= \frac{v(rsab)_{rsab}}
{\epsilon_a + \epsilon_b - \epsilon_r - \epsilon_s},
\label{eq:tamp}\\[4pt]
E^{(20)}_{\mathrm{disp}} &= 4 \sum_{abrs} t_{rsab}\, v(abrs)_{abrs} .
\label{eq:disp}
\end{align}
Under Eq.~\eqref{eq:pairfac} the numerator is
$v(abrs)_{abrs} = \sum_P b^{P}_{(ar)}[a,r]\, b^{P}_{(bs)}[b,s]$, i.e. the
familiar DF-MP2 form with the auxiliary range extended by one.

\paragraph{Which orbital energies (implementation note).} The denominators in
Eq.~\eqref{eq:tamp} use the \emph{CQED} orbital energies. The perturbation
series is defined relative to the CQED monomer Hamiltonians and the numerator
already carries cavity-dressed integrals, so canonical (cavity-free)
denominators would mix two reference definitions. The code retains a
\texttt{canonical\_denom} switch as a \emph{diagnostic} only, and records which
convention produced a given set of amplitudes so that the
exchange-dispersion term cannot silently consume amplitudes built under the
other one.

\subsection{Long-range behaviour}
\label{sec:longrange}

The cavity contribution to Eq.~\eqref{eq:disp} deserves comment because it
looks pathological and is not. Its kernel, $d_{pq}d_{rs}$, contains \emph{no
Coulomb operator} and therefore no $R$-dependence at any separation. Through
$v(abrs)$ it becomes $d^{A}_{ar} d^{B}_{bs}$ in the numerator; the denominator
also tends to a constant as $R\to\infty$, so the cavity part of
$E^{(20)}_{\mathrm{disp}}$ approaches a finite non-zero constant rather than
decaying.

This is the single-mode long-wavelength approximation behaving as specified:
both monomers couple to the same cavity mode with the same $\lam$, so the
cavity mediates an interaction with no distance decay.

The $R\to\infty$ limit is computable in closed form from \emph{isolated-monomer}
transition dipoles and excitation energies, with no dimer calculation:
\begin{equation}
\lim_{R\to\infty} E^{(20)}_{\mathrm{disp}}
= 4 \sum_{ar}\sum_{bs}
\frac{(d^{A}_{ar})^{2}\,(d^{B}_{bs})^{2}}
     {\epsilon_a + \epsilon_b - \epsilon_r - \epsilon_s} .
\label{eq:plateau}
\end{equation}
That makes it a sharp validation target rather than a qualitative observation.
Running it is instructive, and gives a mixed verdict.

\paragraph{The numerator plateaus; the energy does not.} Measured on water/He,
$\lVert v_{\mathrm{cav}}\rVert$ is constant to seven significant figures from
$R = 8$ to $50$\,\AA{}, exactly as the vanishing Coulomb kernel requires. The
\emph{energy}, however, decays steadily, reaching only $5.5\%$ of
Eq.~\eqref{eq:plateau} at $50$\,\AA. With the numerator constant, the drift is
entirely in the denominators.

\paragraph{Implementation note: the orbital energies are not translation
invariant.} The cause is a property of the CQED-SCF reference rather than of
the dispersion expression. For a single atom translated by $z_0$ along the
polarization direction --- physically the same system --- the total CQED-SCF
energy is invariant to $4\times10^{-14}\,E_h$, but the HOMO--LUMO gap grows as
$\lambda^2 z_0^2$, and occupied and virtual orbitals move in \emph{opposite}
directions. Consequently a rigid translation of an entire dimer by $20$\,\AA{}
leaves $E^{(10)}_{\mathrm{elst}}$, $E^{(10)}_{\mathrm{exch}}$,
$E^{(20)}_{\mathrm{ind,r}}$ and $E^{(20)}_{\mathrm{exch-ind,r}}$ unchanged to
$10^{-12}\,E_h$ while changing $E^{(20)}_{\mathrm{disp}}$ by a factor of $7.4$.

Only the terms dividing by \emph{bare} monomer orbital-energy differences are
affected. Induction is invariant despite also involving orbital energies,
because its coupled response carries the compensating DSE Hessian contribution;
Eq.~\eqref{eq:tamp} divides by the bare difference with nothing to cancel it.
The opposite-sign shift of occupied against virtual orbitals points at the
exchange-like term $N_{\mu\nu}[D] = \sum_{\rho\sigma} d_{\mu\rho} d_{\nu\sigma}
D_{\rho\sigma}$ of Eq.~(fock) in the companion document, which is built from
the bare dipole matrix rather than the coherent-state fluctuation
$\dop - \dexp$: under a translation $d \to d + z_0\lambda S$ this acquires a
$z_0^2\lambda^2\,(SDS)$ contribution supported only on the occupied space.

This is independent of the factorization developed in this note --- it is
present identically in the dense four-index formulation, and absent at
$\lam = \bm{0}$ --- but it must be settled before any long-range
cavity-dispersion curve is interpreted physically. Section~\ref{sec:origin}
sets out the three remedies under consideration.

% =============================================================================
\section{Origin dependence of the reference: three remedies}
\label{sec:origin}

\subsection{The defect, stated precisely}

Under a rigid translation by $\bm{T}$, the AO dipole matrix transforms as
$d \to d + (\lam\cdot\bm{T})\,S$. Two terms of the Pauli--Fierz Fock operator
respond:

\begin{itemize}[leftmargin=*, itemsep=3pt]
\item the one-electron DSE (quadrupole) term $Q^{\mathrm{PF}}$ acquires a
contribution proportional to $S$, shifting \emph{every} orbital by
$+\tfrac12(\lam\cdot\bm{T})^2$;
\item the exchange-like term $N_{\mu\nu}[D] = \sum_{\rho\sigma} d_{\mu\rho}
d_{\nu\sigma} D_{\rho\sigma}$ acquires $(\lam\cdot\bm{T})^2 (SDS)$, which is
supported \emph{only on the occupied space}, so $-\bm{N}$ shifts occupied
orbitals by $-(\lam\cdot\bm{T})^2$.
\end{itemize}

Net: occupied $-\tfrac12(\lam\cdot\bm{T})^2$, virtual
$+\tfrac12(\lam\cdot\bm{T})^2$, and every gap widens by $(\lam\cdot\bm{T})^2$.
Measured for water at $T = 20$\,\AA{}, $\lambda = 0.1$: occupied shift
$-7.142242$ against $-\tfrac12\lambda^2T^2 = -7.142130$.

What is \emph{not} broken matters as much. The total CQED-SCF energy is
invariant to $4\times10^{-14}\,E_h$ and the density to $8\times10^{-12}$, so the
occupied and virtual \emph{subspaces} are correct; it is the canonical rotation
within each subspace, and the eigenvalues, that move. Consequently
$E^{(10)}_{\mathrm{elst}}$, $E^{(10)}_{\mathrm{exch}}$,
$E^{(20)}_{\mathrm{ind,r}}$ and $E^{(20)}_{\mathrm{exch-ind,r}}$ are all
invariant; only the terms dividing by bare orbital-energy differences are not.

\subsection{Option 1: an intrinsic reference frame (implemented)}

Solve each monomer's CQED-SCF after translating its ghosted geometry so that
its own real-atom centre of mass lies at the origin. Every monomer quantity is
then independent of where the dimer sits, and translation invariance follows by
construction. The MO coefficients transfer into the dimer frame unchanged,
because a rigid translation moves the basis functions with the molecule.

The implementation turns on one distinction, which is the substantive content
of this option:

\begin{quote}
A block whose four orbital indices all belong to \emph{one} monomer is an
\textbf{internal} quantity and must be built in that monomer's intrinsic frame.
Every block that mixes the monomers is an \textbf{interaction} quantity and must
be built in the shared dimer frame.
\end{quote}

Both halves are load-bearing.

\emph{Interaction quantities.} Monomers $A$ and $B$ use \emph{different}
intrinsic frames, so dipole matrices taken from those frames would differ by
$(\lam\cdot(\bm{T}_A-\bm{T}_B))S$. That would misstate the interaction and
destroy $d^A = d^B$, and with it the single-auxiliary-row factorization
Eq.~\eqref{eq:augmented} on which the entire density-fitted backend rests. The
overlap, the ERIs, the dipole matrices and the nuclear attraction integrals
therefore all stay in the dimer frame.

\emph{Internal quantities.} The coupled-perturbed Hartree--Fock matrix is
monomer $B$'s own orbital Hessian, built from the blocks $v(sbbs)$ and
$v(sbsb)$. Its diagonal carries $\epsilon_s - \epsilon_b$ from the intrinsic
frame, while its exchange-like DSE contribution $-d_{bb'}d_{ss'}$ involves the
occupied--occupied and virtual--virtual blocks of $d$, which are \emph{not}
translation invariant (unlike the occupied--virtual blocks, which are, because
$S_{bs} = 0$). Induction is origin independent only because these two
origin-dependent pieces cancel; supplying intrinsic-frame orbital energies
alongside a dimer-frame Hessian breaks the cancellation. In an intermediate
version of this implementation that mismatch changed $E^{(20)}_{\mathrm{ind,r}}$
by $1.1\times10^{-5}\,E_h$ and reversed its sign under a $20$\,\AA{}
translation. With the Hessian built in the matching frame, induction returns
\emph{exactly} to its dimer-frame value.

\paragraph{Result.} All six components become translation invariant, and the
cavity dispersion plateaus as Eq.~\eqref{eq:plateau} predicts:

\begin{center}
\small
\begin{tabular}{lrr}
\toprule
& \multicolumn{2}{c}{change under a $20$\,\AA{} rigid translation} \\
\cmidrule(l){2-3}
component & \texttt{dimer} frame & \texttt{monomer\_com} frame \\
\midrule
$E^{(10)}_{\mathrm{elst}}$       & $4.7\times10^{-14}$ & $4.5\times10^{-14}$ \\
$E^{(10)}_{\mathrm{exch}}$       & $5.6\times10^{-14}$ & $8.0\times10^{-17}$ \\
$E^{(20)}_{\mathrm{disp}}$       & $\bm{1.8\times10^{-5}}$ & $1.0\times10^{-18}$ \\
$E^{(20)}_{\mathrm{exch-disp}}$  & $\bm{6.0\times10^{-7}}$ & $4.6\times10^{-18}$ \\
$E^{(20)}_{\mathrm{ind,r}}$      & $7.1\times10^{-15}$ & $2.0\times10^{-17}$ \\
$E^{(20)}_{\mathrm{exch-ind,r}}$ & $8.9\times10^{-15}$ & $1.8\times10^{-17}$ \\
\bottomrule
\end{tabular}
\end{center}

\noindent
$\mathrm{Disp20[cav,cav]}$ is then constant at $-2.909131200\times10^{-5}\,E_h$
from $R = 8$ to $50$\,\AA{} (eleven significant figures), against the
isolated-monomer prediction of Eq.~\eqref{eq:plateau},
$-2.9087\times10^{-5}\,E_h$; the residual difference is the ghosted (dimer)
basis, as expected. Only the dispersion terms change value; electrostatics,
exchange and induction are reproduced exactly.

\paragraph{What Option 1 does not do.} It selects a convention rather than
removing one. The residual dependence enters through $\langle z^2\rangle$, so
centre of mass, centre of nuclear charge and centroid of electronic charge give
different answers. Centre of mass is also a \emph{mass}-weighted choice in a
problem with no mass dependence, which is hard to justify from the Hamiltonian.

\subsection{Option 2: build the DSE Fock terms from the fluctuation operator}

The coherent-state DSE is $\tfrac12(\lam\cdot(\hat{\bm\mu}-\langle\hat{\bm\mu}
\rangle))^2$, and the operator inside the square is manifestly translation
invariant. The origin dependence in the Fock arises because the implementation
expands the square and builds $Q^{\mathrm{PF}}$ and $\bm{N}[D]$ from the
\emph{bare} $d$ rather than from the fluctuation. Restoring it costs one
definition:
\begin{equation}
c = \frac{\dexp_{\mathrm{el}}}{N_{\mathrm{el}}},
\qquad
\tilde{d} = d - c\,S ,
\label{eq:fluct}
\end{equation}
where $c \to c + \lam\cdot\bm{T}$ under translation, so $\tilde d$ is exactly
invariant. Using $\tilde d$ in $\bm{N}[D]$ and correcting the one-electron term
to $\bm{H} = \bm{H}_0 + \bm{Q}^{\mathrm{PF}} - (c\,d - \tfrac12 c^2 S)$ --- which
is $-\tfrac12(\tilde d^{\,2})_{\mathrm{AO}}$ expanded --- gives:

\begin{center}
\small
\begin{tabular}{lrrl}
\toprule
Fock construction & $\max|\Delta D|$ & $\max|\Delta\epsilon|$ & HOMO--LUMO gap, $0 \to 20$\,\AA \\
\midrule
bare $d$ (current) & $3.2\times10^{-10}$ & $7.61$ & $0.69545907 \to 15.27376202$ \\
fluctuation $\tilde d$ & $1.8\times10^{-13}$ & $2.6\times10^{-13}$ & $0.69367500 \to 0.69367500$ \\
$\tilde d$, wrong sign & $4.9\times10^{-1}$ & $14.9$ & converges elsewhere \\
\bottomrule
\end{tabular}
\end{center}

\noindent
Advantages over Option 1: no convention is chosen, $d^A = d^B$ is untouched, and
the internal/interaction distinction of Section~\ref{sec:origin} disappears
because there is nothing left to be inconsistent about. Costs: it is a different
definition of the mean-field operator, so it changes results even at the origin
(the gap moves from $0.69545907$ to $0.69367500$) and every pinned reference
value would need re-establishing. The prototype validated the Fock only; whether
the total-energy expression remains consistent under the same substitution has
not been checked, and the sign had to be pinned empirically.

\subsection{Option 3: strong-coupling QED-HF}

Riso, Haugland, Ronca and Koch~\cite{Riso2022} identify origin dependence of the
QED-HF reference as a structural defect rather than an implementation detail,
and resolve it by generalising the ansatz: the coherent-state transformation is
made orbital-dependent and optimised together with the orbitals, rather than
fixed by the mean-field dipole expectation. This removes the origin dependence
of the orbitals and orbital energies themselves, at the level of the reference
definition, and therefore fixes every downstream quantity at once rather than
term by term.

It is the most principled of the three and the most invasive: it replaces the
CQED-SCF reference, so every pinned energy in the package would be
re-established against a different reference, and the SAPT monomer interface
would need to expose whatever the generalised transformation carries.

\subsection{Summary}

\begin{center}
\small
\renewcommand{\arraystretch}{1.25}
\begin{tabular}{lccc}
\toprule
& Option 1 (frame) & Option 2 (fluctuation) & Option 3 (SC-QED-HF) \\
\midrule
removes origin dependence & yes & yes & yes \\
chooses a convention & \textbf{yes} & no & no \\
changes the reference definition & no & \textbf{yes} & \textbf{yes} \\
changes results at the origin & dispersion only & \textbf{all terms} & \textbf{all terms} \\
preserves $d^A = d^B$ & yes (by design) & yes & yes \\
needs internal/interaction split & \textbf{yes} & no & no \\
pinned values need re-establishing & dispersion only & \textbf{all} & \textbf{all} \\
status & implemented, tested & prototyped & not started \\
\bottomrule
\end{tabular}
\end{center}

\noindent
Option 1 is available now as
\texttt{monomer\_reference\_frame="monomer\_com"} and is opt-in; the historical
\texttt{"dimer"} behaviour remains the default so that existing reference values
are untouched. It is best read as a stopgap that makes long-range cavity
dispersion interpretable, not as a resolution of the underlying question, which
Options 2 and 3 address at the level where it arises.

\subsection{Exchange-dispersion}

$E^{(20)}_{\mathrm{exch\text{-}disp}}$ contracts the same amplitudes against an
overlap-dressed intermediate,
\begin{equation}
E^{(20)}_{\mathrm{exch\text{-}disp}}
= -2 \sum_{abrs}\Bigl[\tilde v(absr) + h_{absr} + q_{absr}\Bigr] t_{rsab},
\label{eq:exchdisp}
\end{equation}
where $h$ and $q$ are built from $\tilde v$ blocks contracted with the
intermonomer overlap blocks $s(ab)$, $s(as)$, $s(br)$. The algebra is lengthy
but structurally uniform: every $\tilde v$ block entering
Eq.~\eqref{eq:exchdisp} is one of the nine strings
$abar$, $abra$, $absb$, $abbs$, $abas$, $abrb$, $abab$, $abrs$, $absr$, and
each factorizes by Eq.~\eqref{eq:pairfac}. No $\nbf^4$ AO quantity is required
at any point; the largest intermediates are of dimension
$n_o^A n_o^B n_v^A n_v^B$, which the dense formulation also forms.

% =============================================================================
\section{Numerical verification}
\label{sec:numerics}

The factorization was checked against the dense four-index reference for a
water/He dimer in cc-pVDZ at $\lam = (0,0,0.1)$, over every four-character
string used by the driver (seventeen strings, including all cross-monomer pair
types). The asymmetry between the two rows below is the substantive result:

\begin{center}
\small
\begin{tabular}{lll}
\toprule
context & $\max$ deviation from dense $v$ & interpretation \\
\midrule
\textsf{standard} & $7\times10^{-6}$ -- $6\times10^{-5}$ &
  ordinary fitting error on MO integrals \\
\textsf{cavity} & $3.3\times10^{-17}$ &
  \textbf{exact} --- as Eq.~\eqref{eq:augmented} requires \\
\bottomrule
\end{tabular}
\end{center}

\noindent
Had the DSE row been mispaired, transposed, or misscaled, the cavity row would
read like the standard row. The eleven orders of magnitude between them is what
makes this a decisive rather than a suggestive test.

\paragraph{External oracle.} At $\lam = \bm 0$ the theory reduces exactly to
conventional SAPT0, so Psi4's own \texttt{sapt0} is an independent
implementation against which the dense reference can be pinned. For the same
dimer the dense $E^{(20)}_{\mathrm{disp}}$ and
$E^{(20)}_{\mathrm{exch\text{-}disp}}$ agree with Psi4 to
$3.9\times10^{-8}$ and $2.5\times10^{-8}$ $E_h$ respectively, consistent with
Psi4's density-fitting error against our exact-ERI reference. (Psi4's
\texttt{FISAPT} routines are \emph{functional-group} SAPT and are not the
appropriate oracle here.)

\paragraph{Partition identity.} \textsf{standard} $+$ \textsf{cavity} $=$
\textsf{total} is verified to $10^{-13}$ for every string; by
Eq.~\eqref{eq:contexts} this is a slicing identity, so any deviation would
indicate an indexing error rather than a physical discrepancy.

\paragraph{End-to-end.} With the role assignment above, all six QED-SAPT0
components computed through the factorization agree with the dense four-index
reference for the water/methylamine dimer in jun-cc-pVDZ at
$\lam = (0,0,0.1)$:

\begin{center}
\small
\begin{tabular}{lrrr}
\toprule
component & dense / $E_h$ & factorized / $E_h$ & difference / kcal\,mol$^{-1}$ \\
\midrule
$E^{(10)}_{\mathrm{elst}}$      & $-0.0086737093$ & $-0.0086743054$ & $-3.7\times10^{-4}$ \\
$E^{(10)}_{\mathrm{exch}}$      & $ 0.0030805286$ & $ 0.0030819372$ & $ 8.8\times10^{-4}$ \\
$E^{(20)}_{\mathrm{disp}}$      & $-0.0027151107$ & $-0.0027146006$ & $ 3.2\times10^{-4}$ \\
$E^{(20)}_{\mathrm{exch-disp}}$ & $ 0.0001966590$ & $ 0.0001952441$ & $-8.9\times10^{-4}$ \\
$E^{(20)}_{\mathrm{ind,r}}$     & $-0.0017887767$ & $-0.0017887293$ & $ 3.0\times10^{-5}$ \\
$E^{(20)}_{\mathrm{exch-ind,r}}$& $ 0.0008167150$ & $ 0.0008166473$ & $-4.2\times10^{-5}$ \\
\midrule
total                           & $-0.0090836941$ & $-0.0090838066$ & $-7.1\times10^{-5}$ \\
\bottomrule
\end{tabular}
\end{center}

\noindent
with $N_{\mathrm{bf}} = 89$, $\naux = 398$ (JKFIT) and $293$ (RIFIT). Wall time
for the integral construction and all six components fell from $11.7$\,s to
$1.1$\,s and peak memory from $2380$\,MB to $254$\,MB.

% =============================================================================
\section{Relationship to the treatment of the DSE elsewhere in the package}

The factorization developed here is a restatement, not a new approximation, of
the position taken throughout \texttt{cqed-scf}: the DSE is a genuine
two-electron operator. This is worth recording explicitly because the
literature is not uniform on the point. Yang \emph{et al.}~\cite{Yang2021}
evaluate the DSE as the square of a one-body expectation value over a single
determinant, which yields a rank-one \emph{mean-field} coupling with no
exchange-like counterpart. Treating it as a two-electron operator instead
produces the exchange-like terms that appear as $-N_{\mu\nu}[D]$ in the
Pauli--Fierz Fock matrix, as $K_{\mathrm{DSE}}$ in the \texttt{DSEJK} builder,
and as $-d_{ij}d_{ab}$ in the response block. These are different levels of
theory, not a rescaling of one another.

The three faces of the same operator now implemented in the package are:

\begin{center}
\small
\begin{tabular}{lll}
\toprule
face & object & used by \\
\midrule
$J/K$ build & $J = d\,\Tr(dD)$, \; $K = d D d^{\mathsf T}$ & SCF, SAPT first order, induction \\
MO response block & $2 d_{ia} d_{jb} - d_{ij} d_{ab}$ & QED-CIS, CPHF \\
DF auxiliary row & $\tilde B^{\naux+1} = d$ & SAPT dispersion (this note) \\
\bottomrule
\end{tabular}
\end{center}

\noindent
All three follow from Eq.~\eqref{eq:sepkernel} by the substitution
$(pq|rs)\to d_{pq}d_{rs}$ in the corresponding standard expression. The
response face evaluates it directly in the MO blocks (cheapest: no AO
transform and no $J/K$ build at all); the SAPT face needs the AO form.

% =============================================================================
\begin{thebibliography}{9}

\bibitem{Haugland2020}
T.~S. Haugland, E.~Ronca, E.~F. Kj{\o}nstad, A.~Rubio, and H.~Koch,
\emph{Coupled Cluster Theory for Molecular Polaritons: Changing Ground and
Excited States}, Phys. Rev. X \textbf{10}, 041043 (2020).

\bibitem{Yang2021}
J.~Yang, Q.~Ou, Z.~Pei, H.~Wang, B.~Weng, Z.~Shuai, K.~Mullen, and Y.~Shao,
\emph{Quantum-electrodynamical time-dependent density functional theory within
Gaussian atomic basis}, J. Chem. Phys. \textbf{155}, 064107 (2021).

\bibitem{Jeziorski1994}
B.~Jeziorski, R.~Moszynski, and K.~Szalewicz,
\emph{Perturbation Theory Approach to Intermolecular Potential Energy Surfaces
of van der Waals Complexes}, Chem. Rev. \textbf{94}, 1887 (1994).

\bibitem{Hohenstein2012}
E.~G. Hohenstein and C.~D. Sherrill,
\emph{Raising the Prediction Limit: Symmetry-Adapted Perturbation Theory},
WIREs Comput. Mol. Sci. \textbf{2}, 304 (2012).

\bibitem{Riso2022}
R.~R. Riso, T.~S. Haugland, E.~Ronca, and H.~Koch,
\emph{Molecular orbital theory in cavity QED environments},
Nat. Commun. \textbf{13}, 1368 (2022).

\bibitem{Whitten1973}
J.~L. Whitten,
\emph{Coulombic potential energy integrals and approximations},
J. Chem. Phys. \textbf{58}, 4496 (1973).

\bibitem{Dunlap1979}
B.~I. Dunlap, J.~W.~D. Connolly, and J.~R. Sabin,
\emph{On some approximations in applications of $X\alpha$ theory},
J. Chem. Phys. \textbf{71}, 3396 (1979).

\end{thebibliography}

\end{document}
